\scanpage{30}%
\mspage{15}%
\[
y^2=x^3-x \qquad (\text{si }c_4=c_6=0) \tag*{(6.7)}
\]

c) $p=2$ : ici, on met la courbe sour la forme (6.3) et on résoud (6.4) avec $a_3\unsure{\cdot}=1$, $a_4=a_6=0$ pour obtenir la forme type
\[
y^2+y=x^3 \qquad (\text{si }c_4=c_6=0) \, . \tag*{(6.8)}
\]

Pour compléter le tableau, reste à décrire les automorphismes des courbes elliptiques.

\uline{Proposition} 6.9. \quad \uline{Soit} $E$ \uline{une courbe de genre un sur une base} $S$, \uline{sans fibre additive}.

(I) \quad \uline{Si} $c_4$ et $c_6$ \uline{sont inversibles,} \uline{les seuls automorphismes}\margin{\tikz[baseline=(rem.base)]{\node[align=left,inner sep=2pt] (rem) {\handwritten{faux si $S$}\\\handwritten{\unsure{pas} connexe}};\handwritten{\draw[thin,gray,rounded corners=4pt] (rem.north east) -- (rem.south east) -- (rem.south west) -- (rem.north west) -- ([xshift=-0.4em]rem.north east);\draw[thin] ([xshift=0.35em,yshift=0.5em]rem.north west) -- ([xshift=0.2em]rem.west);}}} \uline{de} $E$ \uline{sont l'identité et la symétrie}.

(II) \quad \uline{Si} $2$ et $3$ \uline{sont inversibles}, \uline{l'application canonique}
\diagram{\[
\vcenter{\hbox{\begin{tikzcd}[ampersand replacement=\&,column sep=1.9em]
 \text{\uline{Aut}}(E) \handwritten{\arrow[r]} \& G_m
\end{tikzcd}}}\tag*{(6.6.)}
\]}
\uline{donnée par l'action sur les différentielles invariantes}, \uline{est une immersion fermée}.

(III) \uline{Si} $2$ \uline{et} $3$ \uline{sont inversibles}, \uline{et} $c_4=0$ (resp. $c_6=0$) \uline{l'image de} (6.6) \uline{est} $\mu_6$ (resp. $\mu_4$).

\uline{Supposons maintenant que} $S$ \uline{soit le spectre d'ur corps}\tikz[baseline=(stop.base)]{\node[inner sep=0pt] (stop) {.};\handwritten{\draw[thin,gray,rotate around={-18:(stop.center)}] (stop.center) ellipse (0.57em and 1.03em);}}

(IV) \quad \uline{Si} $3=c_4=c_6=0$, $\mathrm{Aut}(E)$ \uline{est géométriquement \unsure{isomorphe} au produit semi-direct} (\uline{non commutatif}) \uline{de} $\mathbb{Z}/(4)$ \uline{par} $\mathbb{Z}/(3)$, \uline{et l'image de} (6.6) \uline{est} $\mu_4$ .

(V) \quad \uline{Si} $2=c_4=c_6=0$, $\mathrm{Aut}(E)$ \uline{est géométriquement d'ordre} \tikz[baseline=(order.base)]{\node[inner sep=0pt] (order) {$24$,};\handwritten{\draw[thin,gray,rotate around={-38:(order.center)}] ([xshift=.37em,yshift=2.08em]order.center) arc[start angle=72,end angle=365,x radius=1.19em,y radius=2.2em];}}\margin{\handwritten{\tikz[baseline=-0.3ex]{\draw[thin] (0,0.7em) -- (0.35em,-0.05em);} structure ?}} \uline{et l'image de} (6.6) \uline{est} $\mu_3$.
% archive sheet number: 29
